\section{Limitations}
\label{sec:limitations}

We are explicit about what we do and do not establish.

\paragraph{Node-level, not pod-level.}
On confidential-VM node pools the attestation evidence refers to the confidential
worker node/VM, not to a per-pod TEE. We therefore report \emph{attested-node
placement}; we do not claim pod-level workload attestation, which would require a
confidential-container runtime (CoCo/Kata-CC) enabled and measured.

\paragraph{No confidential GPU, no TDX.}
The confidential hardware used is CPU-only AMD SEV-SNP (\texttt{DC8as\_v6}). We do
not evaluate confidential GPUs (e.g.\ NVIDIA H100 CC) or Intel TDX; GPU-aware
placement is expressed in the policy schema but its validation is future work.

\paragraph{AI workload scope.}
The AKS evaluation includes real governed AI workloads (OpenAI chat, OpenAI
embedding, and a local CPU model) to validate attested placement, model/image
digest binding, and offline placement verification. These workloads do not
establish model confidentiality, OpenAI service-side confidentiality,
confidential GPU execution, or large-model throughput.

\paragraph{Latency scope.}
The AKS timing claim is based on a B1--B5 campaign with 30 measured runs per
baseline. The B1--B5 figure uses one comparable client-observed timing
instrument; under that instrument, B5's median is 1167.0\,ms. B5 also records
an in-process scheduler-internal phase metric (median 216.232\,ms), but that
metric is reported separately and must not be read as an end-to-end speedup over
B1--B4. We report prototype overhead, not a universal Kubernetes latency bound.

\paragraph{Scale.}
The final AKS campaign scaled the \texttt{Standard\_DC8as\_v6} confidential pool
to four active nodes under a 32-vCPU DCasv6 quota. This supports multi-node node
qualification and B1--B5 overhead measurement, but it is not a cloud-scale
throughput study. Local \texttt{kind} and KWOK harnesses are retained only for
CI/debug/regression or scheduler-only scalability context and are not presented
as main cloud security or performance results.

\paragraph{MAA as trust anchor.}
Real evidence trusts the Microsoft Azure Attestation service and the SEV-SNP
hardware within their stated threat model. A verifier that pins MAA signing keys
and policy is used; a fully self-hosted verifier is out of scope.

\paragraph{Intellectual property.}
The placement-token and signed-decision mechanisms are subject to an IP review
before submission (marked in the design/implementation sections); the artifact is
private pending that review.
